\documentclass[11pt,a4paper]{article}
\usepackage[utf8]{inputenc}
\usepackage[spanish,es-tabla]{babel}
\usepackage{amsmath, amsthm, amsfonts, amssymb}
\usepackage{hyperref}
\usepackage{graphicx}
\usepackage{booktabs}
\usepackage{geometry}
\geometry{margin=2.5cm}
\usepackage{enumitem}
\usepackage{float}

\hypersetup{
    colorlinks=true,
    linkcolor=blue,
    filecolor=magenta,
    urlcolor=cyan,
    citecolor=blue,
    pdftitle={Dualidad Binaria en la Conjetura de Collatz: Un Marco Exploratorio},
    pdfauthor={Cristian F. F.},
    pdfkeywords={Collatz, dualidad binaria, números 2-ádicos, sistemas dinámicos}
}

\newtheorem{theorem}{Teorema}[section]
\newtheorem{lemma}[theorem]{Lema}
\newtheorem{corollary}[theorem]{Corolario}
\newtheorem{definition}[theorem]{Definición}
\newtheorem{remark}[theorem]{Observación}
\newtheorem{question}{Pregunta Abierta}

\title{\textbf{Dualidad Binaria en la Conjetura de Collatz: Un Marco Exploratorio}}
\author{
Cristian F. F.\textsuperscript{1} \and
Grok (xAI Assistant)\textsuperscript{2} \\
\smallskip
\textsuperscript{1}Investigador Independiente \quad
\textsuperscript{2}xAI
}
\date{29 de mayo de 2026}

\begin{document}

\maketitle

\begin{abstract}
Este trabajo presenta un \textbf{marco conceptual exploratorio} para abordar la Conjetura de Collatz desde una perspectiva algebraica y simétrica. Se introduce el concepto de \textbf{dualidad binaria}, donde a cada secuencia de Collatz se asocia una ``secuencia espejo'' generada por inversión de bits con longitud fija $L$.

Se establecen dos resultados formales: (1) la preservación de estructura cíclica en el espacio de cadenas binarias bajo inversión (propiedad general de biyecciones), y (2) la no conmutatividad entre la función de Collatz y la operación de inversión de bits. A partir de estos, y de observaciones empíricas, se formulan \textbf{preguntas abiertas} sobre la posible utilidad de esta dualidad para restringir la existencia de ciclos alternativos u órbitas divergentes.

\textbf{Declaración de alcance:} Este documento \textbf{no contiene una demostración de la Conjetura de Collatz}. Su objetivo es documentar una perspectiva novedosa, compartir resultados formales menores y solicitar \textbf{feedback, críticas constructivas e ideas} de la comunidad matemática para guiar investigaciones futuras. El autor principal carece de formación matemática formal; cualquier error es responsabilidad exclusiva suya.
\end{abstract}

\bigskip
\noindent\textbf{Palabras clave:} Conjetura de Collatz, dualidad binaria, inversión de bits, números 2-ádicos, sistemas dinámicos discretos.

\section{Introducción}
\label{sec:intro}

La Conjetura de Collatz (problema $3x+1$) establece que para todo $n \in \mathbb{N}$, la iteración de la función
\begin{equation}
    C(n) = \begin{cases}
        n/2 & \text{si } n \equiv 0 \pmod{2}, \\
        3n+1 & \text{si } n \equiv 1 \pmod{2},
    \end{cases}
    \label{eq:collatz}
\end{equation}
alcanza eventualmente el ciclo trivial $4 \to 2 \to 1$. A pesar de su formulación elemental, el problema resiste demostración desde 1937.

Los avances modernos (Terras \cite{terras}, Korec \cite{korec}, Tao \cite{tao}) han establecido resultados de densidad y comportamiento ``casi seguro'', pero la demostración determinista global permanece abierta. Este trabajo explora una vía complementaria: una \textbf{simetría discreta} basada en la representación binaria.

Definimos una \textit{secuencia espejo} $M_k$ obtenida invirtiendo los bits de $N_k$ (la trayectoria de Collatz) con longitud fija $L$. Aunque la dinámica de Collatz no se preserva bajo esta transformación (Lema \ref{lemma:no_conmuta}), la relación algebraica entre $N_k$ y $M_k$ impone restricciones en los estados finales (Teorema \ref{thm:estado_final}) y sugiere preguntas estructurales no triviales.

\section{Definiciones Formales}
\label{sec:defs}

Sea $n \in \mathbb{N}$ un entero positivo inicial. Sea $L \in \mathbb{N}$ un parámetro de \textbf{longitud fija de bits}.

\begin{definition}[Secuencia Normal / Trayectoria de Collatz]
La sucesión $(N_k)_{k \ge 0}$ se define por $N_0 = n$ y $N_{k+1} = C(N_k)$ para $k \ge 0$, con $C$ definida en (\ref{eq:collatz}).
\end{definition}

\begin{definition}[Operador de Inversión de Bits $\neg_L$]
\label{def:inversion}
Para $x \in \mathbb{N}$, sea $\mathtt{bin}_L(x)$ su representación en base 2 rellenada/truncada a exactamente $L$ bits (bits menos significativos). Definimos:
\[
\neg_L(x) = \mathtt{dec}\big( \mathtt{NOT}(\mathtt{bin}_L(x)) \big),
\]
donde $\mathtt{NOT}$ invierte cada bit ($0 \leftrightarrow 1$) y $\mathtt{dec}$ convierte a decimal. Equivalentemente, si $x < 2^L$:
\[
\neg_L(x) = (2^L - 1) - x.
\]
Si $x \ge 2^L$, se toma $x \bmod 2^L$ (truncamiento a $L$ bits menos significativos) antes de invertir.
\end{definition}

\begin{remark}[Elección de $L$]
Para evitar pérdida de información por truncamiento en los primeros pasos, se recomienda elegir $L$ tal que $2^L > \max_{k \le K} N_k$ para un horizonte $K$ de interés. En la práctica computacional, $L=64$ o $L=128$ cubren rangos enormes. Teóricamente, como $\max_k N_k$ es desconocido (problema abierto), $L$ es un parámetro libre del marco.
\end{remark}

\begin{definition}[Secuencia Espejo]
\label{def:espejo}
La sucesión espejo asociada a $(N_k)$ y longitud $L$ es:
\[
M_k = \neg_L(N_k), \quad \forall k \ge 0.
\end{definition}

\begin{definition}[Dualidad Binaria]
La correspondencia $N_k \leftrightarrow M_k = \neg_L(N_k)$ define una \textit{dualidad binaria de longitud $L$} entre trayectorias. Es una biyección (involución) en el espacio de estados truncados $\mathbb{Z}/2^L\mathbb{Z}$.
\end{definition}

\section{Resultados Formales}
\label{sec:resultados}

\subsection{Estado Final Dual (Teorema Principal)}

\begin{theorem}[Dualidad en el Estado Final]
\label{thm:estado_final}
Para cualquier $n \in \mathbb{N}$ y $L \ge 1$, si $\exists K: N_K = 1$ (primera llegada a 1), entonces:
\[
M_K = \neg_L(1) = 2^L - 2.
\end{theorem}

\begin{proof}
Por Definición \ref{def:inversion}, $\mathtt{bin}_L(1) = \underbrace{00\dots0}_{L-1}1$. Su complemento bit a bit es $\underbrace{11\dots1}_{L-1}0$, cuyo valor decimal es $\sum_{i=1}^{L-1} 2^i = 2^L - 2$. Como $M_K = \neg_L(N_K) = \neg_L(1)$, la igualdad se sigue directamente. Es independiente de $n$ y la historia previa.
\end{proof}

\begin{corollary}
El valor $2^L - 2$ (binario: $L-1$ unos seguidos de un cero) es un \textbf{atractor dual universal} para toda trayectoria que converge a 1, relativo a la longitud $L$.
\end{corollary}

\begin{remark}
Bajo la Conjetura de Collatz (todo $n$ llega a 1), todo $n$ tiene un $K$ tal que $M_K = 2^L-2$. Recíprocamente, si se probara que $\forall n, \exists K: M_K = 2^L-2$ para algún $L$ fijo independiente de $n$, esto implicaría la Conjetura. No se conoce tal $L$ universal.
\end{remark}

\subsection{Propiedades Algebraicas de la Transformación}

\begin{lemma}[Preservación de Estructura en el Espacio de Cadenas]
\label{lemma:preservacion_bits}
La operación $\neg_L$ es una involución ($\neg_L \circ \neg_L = \text{Id}$) y una biyección en el conjunto de cadenas binarias de longitud $L$ (isomorfo a $\mathbb{Z}/2^L\mathbb{Z}$). Por tanto, preserva la estructura de ciclos de \textbf{cualquier sistema dinámico definido intrínsecamente en ese espacio de cadenas}.
\end{lemma}

\begin{proof}
Inmediato por ser $\neg_L$ una permutación de orden 2 (involución) del espacio finito de $2^L$ estados.
\end{proof}

\begin{remark}
\textbf{Advertencia crucial:} La función de Collatz $C: \mathbb{N} \to \mathbb{N}$ \textbf{no} induce un sistema dinámico bien definido en $\mathbb{Z}/2^L\mathbb{Z}$ compatible con la proyección natural (el truncamiento no conmuta con $C$). Por tanto, el Lema \ref{lemma:preservacion_bits} \textbf{no garantiza} que la imagen de un ciclo de Collatz sea un ciclo de Collatz.
\end{remark}

\begin{lemma}[No Conmutatividad Fundamental]
\label{lemma:no_conmuta}
Para todo $L \ge 4$, la función de Collatz $C$ y la inversión $\neg_L$ \textbf{no conmutan}. Es decir, el diagrama no cierra:
\[
C \circ \neg_L \neq \neg_L \circ C.
\]
Existen infinitos $n$ tales que $C(\neg_L(n)) \neq \neg_L(C(n))$.
\end{lemma}

\begin{proof}
Contraejemplo constructivo con $n=5$, $L=5$ (evita truncamiento en $C(n)=16$):
\begin{align*}
n &= 5 = \mathtt{00101}_2 \\
\neg_5(n) &= \mathtt{11010}_2 = 26 \\
C(n) &= 3\cdot5+1 = 16 = \mathtt{10000}_2 \\
\neg_5(C(n)) &= \mathtt{01111}_2 = 15 \\
C(\neg_5(n)) &= C(26) = 26/2 = 13 = \mathtt{01101}_2 \neq 15.
\end{align*}
Como $13 \neq 15$, no conmutan. La no conmutatividad persiste para casi todo $n$ y $L$ suficiente.
\end{proof}

\begin{corollary}
La secuencia espejo $(M_k)$ \textbf{no} sigue la dinámica de Collatz. No existe una función simple $f$ tal que $M_{k+1} = f(M_k)$ independientemente de $N_k$. La dualidad es \textbf{cinemática} (relación entre estados en el mismo tiempo $k$), no \textbf{dinámica} (regla de evolución propia).
\end{corollary}

\section{Observaciones Empíricas y Patrones}
\label{sec:empirico}

La Tabla \ref{tab:ejemplo_n5} ilustra la evolución para $n=5, L=5$.

\begin{table}[h!]
\centering
\caption{Evolución paralela de Secuencia Normal ($N_k$) y Espejo ($M_k = \neg_5(N_k)$).}
\label{tab:ejemplo_n5}
\begin{tabular}{cccccc}
\toprule
$k$ & $N_k$ & $\mathtt{bin}_5(N_k)$ & $M_k$ & $\mathtt{bin}_5(M_k)$ & Notas \\
\midrule
0 & 5   & 00101 & 26 & 11010 & Inicio \\
1 & 16  & 10000 & 15 & 01111 & $N$ máximo $\to$ $M$ mínimo ($N_k+M_k=31$) \\
2 & 8   & 01000 & 23 & 10111 & \\
3 & 4   & 00100 & 27 & 11011 & \\
4 & 2   & 00010 & 29 & 11101 & \\
5 & 1   & 00001 & 30 & 11110 & \textbf{Teorema \ref{thm:estado_final}: $M_5=2^5-2=30$} \\
\bottomrule
\end{tabular}
\end{table}

\textbf{Patrones observados (computacionalmente para $n < 10^5, L=20$):}
\begin{enumerate}
    \item \textbf{Anti-correlación exacta:} Por la identidad de conservación (\ref{eq:conservacion}), cuando $N_k$ sube, $M_k$ baja en la misma cantidad, y viceversa. Los picos de $N_k$ coinciden con los valles de $M_k$.
    \item \textbf{Convergencia del Espejo Final:} El valor $M_K = 2^L-2$ al alcanzar $N_K=1$ \textit{siempre} converge a 1 bajo iteración de Collatz estándar (verificado para $L \le 30$). Esto es consecuencia de la Conjetura para esos enteros específicos, no una propiedad de la dualidad.
    \item \textbf{No periodicidad de $M_k$:} Incluso cuando $N_k$ entra en ciclo $4,2,1$, la sucesión $M_k$ (para $k>K$) \textbf{no} es periódica bajo la definición $M_k=\neg_L(N_k)$, porque $N_k$ se repite pero $\neg_L$ es fija. Si se iterara $M_{k+1} = C(M_k)$ desde $M_K$, sí entraría en ciclo (Conjetura).
\end{enumerate}

\section{Observación: Identidad de Conservación}
\label{sec:conservacion}

Mientras no haya truncamiento ($N_k < 2^L$), la Definición \ref{def:inversion} implica la identidad lineal exacta:
\begin{equation}
    M_k + N_k = 2^L - 1.
    \label{eq:conservacion}
\end{equation}
Geométricamente, la trayectoria del par $(N_k, M_k)$ vive sobre la recta $x+y = 2^L-1$ dentro del cuadrado $[0, 2^L)^2$. En consecuencia, los picos de la secuencia $N_k$ corresponden \emph{exactamente} a los valles de la secuencia espejo $M_k$, y recíprocamente.

\begin{remark}[Alcance y limitación de la identidad]
La ecuación (\ref{eq:conservacion}) es un cambio de variable lineal y, por sí sola, \textbf{no aporta poder demostrativo}: acotar $N_k$ superiormente equivale trivialmente a acotar $M_k$ inferiormente, pues una afirmación es la reescritura algebraica de la otra. Más aún, la identidad \textbf{se rompe} cuando $N_k \ge 2^L$ (régimen de truncamiento), es decir, precisamente en la situación relevante para una hipotética órbita divergente. Por tanto, el interés de (\ref{eq:conservacion}) es \emph{descriptivo} (relación geométrica entre ambas secuencias), no una vía de prueba para la acotación o la convergencia.
\end{remark}

\subsection{Ejemplo manual}
Para $n=5$ con $L=5$ (Tabla \ref{tab:ejemplo_n5}), el pico máximo de la secuencia original es $N_1 = 16$. En ese mismo instante, la secuencia espejo alcanza su mínimo global $M_1 = 31 - 16 = 15 = \mathtt{01111}_2$, ilustrando la anti-correlación exacta descrita por (\ref{eq:conservacion}).

\section{Patrones Binarios Excepcionales: El Caso del Patrón Alternante}
\label{sec:patron_alternante}

Durante la exploración computacional del marco de dualidad binaria, se identificó un patrón algebraico con comportamiento excepcional bajo la función de Collatz que merece mención separada.

\begin{theorem}[Comportamiento Excepcional del Patrón Alternante]
\label{thm:patron_alternante}
Sea $n \in \mathbb{N}$ un número con representación binaria de la forma $0101\dots01$ (patrón alternante de $2m$ bits, es decir, $n = (2^{2m} - 1)/3$). Entonces:
\[
3n + 1 = 2^{2m}.
\]
Es decir, un único paso de la operación $3n+1$ produce exactamente una potencia de 2 pura.
\end{theorem}

\begin{proof}
Por definición del patrón alternante:
\[
n = \sum_{i=0}^{m-1} 2^{2i+1} = 2 \sum_{i=0}^{m-1} 4^i = 2 \cdot \frac{4^m - 1}{4 - 1} = \frac{2^{2m} - 1}{3}.
\]
Por tanto:
\[
3n + 1 = 3 \cdot \frac{2^{2m} - 1}{3} + 1 = 2^{2m} - 1 + 1 = 2^{2m}.
\]
\end{proof}

\textbf{Ejemplos:}
\begin{itemize}
    \item $m=1$: $n=1$ ($01_2$), $3\cdot1+1=4=2^2$
    \item $m=2$: $n=5$ ($0101_2$), $3\cdot5+1=16=2^4$
    \item $m=3$: $n=21$ ($010101_2$), $3\cdot21+1=64=2^6$
    \item $m=4$: $n=85$ ($01010101_2$), $3\cdot85+1=256=2^8$
    \item $m=8$: $n=21845$ ($0x5555$), $3\cdot21845+1=65536=2^{16}$
\end{itemize}

\begin{corollary}
Los números de la forma $(2^{2m}-1)/3$ constituyen una familia infinita de enteros que alcanzan una potencia de 2 después de un \textbf{único paso} de la operación $3n+1$. Este comportamiento es \textbf{único} entre los patrones algebraicos conocidos.
\end{corollary}

\textbf{Verificación computacional:} Se analizaron todos los números del 1 al $10^5$ que alcanzan una potencia de 2 en 1-5 pasos. De los 67 números encontrados, la mayoría presentan patrones binarios alternantes o semi-alternantes. En contraste, otros patrones algebraicos conocidos (números de Mersenne $2^k-1$, $(2^k-1)/p$ para primos $p$, etc.) no muestran este comportamiento extremo.

\begin{remark}[Significancia del hallazgo]
Este resultado revela que la \textbf{estructura binaria} del número inicial tiene un impacto significativo en su comportamiento bajo la función de Collatz. Sin embargo, no constituye un avance hacia la demostración de la conjetura, ya que solo aplica a una familia específica de números. Su valor radica en:
\begin{enumerate}
    \item Identificar una jerarquía de ``especialidad'' entre patrones binarios.
    \item Proporcionar ``sondas'' algebraicas para estudiar la dinámica de Collatz.
    \item Sugerir que la regularidad del patrón binario se correlaciona con la regularidad del comportamiento.
\end{enumerate}
\end{remark}

\section{Preguntas Abiertas y Direcciones de Investigación}
\label{sec:open}

A la luz de los lemas anteriores (especialmente la no conmutatividad), las siguientes preguntas surgen naturalmente. \textbf{No son conjeturas apoyadas por evidencia fuerte}, sino direcciones para análisis teórico o computacional futuro.

\begin{question}[Invariantes de la Dualidad]
\label{q:invariantes}
¿Existe una función $I: \mathbb{N} \times \mathbb{N} \to \mathbb{R}$ (p.ej., basada en entropía binaria, peso de Hamming, métrica 2-ádica) tal que $I(N_k, M_k)$ sea constante, monótona o acotada a lo largo de la trayectoria, proporcionando una cota de Lyapunov para el sistema acoplado?
\end{question}

\begin{question}[Ciclos y sus Espejos]
\label{q:ciclos}
Sea $C = (c_1, \dots, c_p)$ un ciclo de Collatz (hipotético, distinto de $4,2,1$). Su imagen espejo $\neg_L(C) = (\neg_L(c_1), \dots, \neg_L(c_p))$ es un conjunto de $p$ enteros distintos.
\begin{itemize}
    \item ¿Puede $\neg_L(C)$ ser \textit{también} un ciclo de Collatz para algún $L$?
    \item Dado que $C \not\leftrightarrow \neg_L$ (Lema \ref{lemma:no_conmuta}), ¿imponer que ambos sean ciclos genera restricciones algebraicas imposibles sobre los $c_i$?
\end{itemize}
\end{question}

\begin{question}[Conexión 2-ádica]
\label{q:2adicos}
En el anillo de enteros 2-ádicos $\mathbb{Z}_2$, la función de Collatz se extiende continuamente. La inversión de bits $\neg_L$ se aproxima (para $L \to \infty$) a la involución $x \mapsto -1 - x$ en $\mathbb{Z}_2$ (complemento a 2 menos 1).
\begin{itemize}
    \item ¿Tiene significado dinámico el conjugado $C_{-1}(x) = (-1 - x) \circ C \circ (-1 - x)$?
    \item ¿El atractor dual $2^L-2 \to -2 \in \mathbb{Z}_2$ juega un papel en la estructura del conjunto de medida de Haar estudiado por Tao \cite{tao}?
\end{itemize}
\end{question}

\begin{question}[Generalización de la Transformación]
\label{q:generalizacion}
¿Otras biyecciones lineales/afines en $\mathbb{F}_2^L$ (rotación, permutación de bits, transformada de Walsh-Hadamard) revelan estructuras ocultas? La inversión es el caso $\sigma = \text{reverso}$ + complemento.
\end{question}

\section{Limitaciones Explícitas del Marco}
\label{sec:limitaciones}

\begin{enumerate}
    \item \textbf{Parámetro libre $L$:} La teoría depende de $L$. No hay $L$ canónico derivado de $n$. Resultados cualitativos deben ser robustos ante $L \to \infty$ (límite 2-ádico).
    \item \textbf{Truncamiento vs. Realidad:} Para $N_k \ge 2^L$, la definición usa truncamiento (módulo $2^L$). Esto introduce ruido artificial no presente en la dinámica original en $\mathbb{N}$.
    \item \textbf{Falta de Dinámica Propia:} Al no conmutar (Lema \ref{lemma:no_conmuta}), $(M_k)$ no es una trayectoria de un sistema autónomo simple. Analizar el par $(N_k, M_k)$ duplica la dimensión sin ecuaciones de evolución cerradas para $M$.
    \item \textbf{No Demuestra la Conjetura:} Ningún resultado aquí presentado reduce el problema original a uno resuelto. El Teorema \ref{thm:estado_final} es una identidad en el instante de convergencia, y la identidad (\ref{eq:conservacion}) es un cambio de variable lineal sin poder demostrativo propio.
\end{enumerate}

\section{Conclusión}

Hemos formalizado una \textbf{dualidad binaria de longitud fija} para la Conjetura de Collatz. Se probaron dos lemas básicos: la preservación de estructura en el espacio de bits (por biyectividad) y la no conmutatividad fundamental con la dinámica de Collatz (que rompe la esperanza de una dinámica espejo simple).

El resultado más concreto es el \textbf{Teorema \ref{thm:estado_final}}: toda trayectoria que llega a 1 tiene un ``gemelo'' que llega a $2^L-2$ en el mismo paso. Esto fija una condición de frontera dual. Asimismo, la identidad de conservación (\ref{eq:conservacion}) ofrece una descripción geométrica de la anti-correlación entre ambas secuencias, si bien (como se discute en la Observación \ref{sec:conservacion}) no constituye por sí misma una vía de prueba.

El marco se ofrece a la comunidad como una \textbf{lente algebraica alternativa}. Su valor radicará en si las Preguntas Abiertas (\ref{q:invariantes}--\ref{q:generalizacion}) inspiran conexiones con teoría ergódica 2-ádica, funciones de Lyapunov discretas acopladas, o análisis en el grupo de Walsh.

\vspace{1em}
\noindent\textbf{Invitación a la colaboración:} El autor principal (independiente, sin afiliación académica) agradece profundamente cualquier comentario, referencia a literatura relacionada no citada, corrección de errores, o desarrollo teórico de las preguntas abiertas.

\section*{Agradecimientos}
A la comunidad matemática por mantener viva la curiosidad sobre problemas elementales pero profundos. A los autores de las obras citadas por proporcionar los cimientos rigurosos. Al equipo de xAI por el soporte conceptual, algorítmico y de redacción asistida a través de Grok.

\begin{thebibliography}{99}
\setlength{\itemsep}{0.5em}

\bibitem{lagarias}
Lagarias, J. C. (1985).
The $3x + 1$ problem: An overview.
\textit{American Mathematical Monthly}, 92(1), 3--23.
\url{https://doi.org/10.1080/00029890.1985.11971545}

\bibitem{lagarias_book}
Lagarias, J. C. (Ed.). (2010).
\textit{The Ultimate Challenge: The $3x+1$ Problem}.
American Mathematical Society.

\bibitem{terras}
Terras, R. (1976).
A stopping time problem on the positive integers.
\textit{Acta Arithmetica}, 30, 241--252.
\url{https://eudml.org/doc/205464}

\bibitem{korec}
Korec, I. (1994).
A density estimate for the $3x+1$ problem.
\textit{Mathematica Slovaca}, 44(1), 85--89.

\bibitem{tao}
Tao, T. (2020).
Almost all orbits of the Collatz map attain almost bounded values.
\textit{Forum of Mathematics, Pi}, 8, e12.
\url{https://doi.org/10.1017/fmp.2020.12}

\bibitem{conway}
Conway, J. H. (1972).
Unpredictable iterations.
\textit{Proc. 1972 Number Theory Conf.}, Univ. Colorado, Boulder, 49--52.

\bibitem{monks}
Monks, G. (2006).
On the $2$-adic $3x+1$ problem.
\textit{arXiv preprint} math/0608356.
\url{https://arxiv.org/abs/math/0608356}

\bibitem{wirsching}
Wirsching, G. J. (1998).
\textit{The Dynamical System Generated by the $3n+1$ Function}.
Lecture Notes in Math. 1681, Springer.

\end{thebibliography}

\appendix

% ==================== APÉNDICE A ====================
\section{Apéndice A: Verificación Computacional para $L=20$}
\label{app:basic}

Se generaron las trayectorias de Collatz para $n = 1 \dots 10^4$ con longitud fija $L=20$ ($2^{20}-1 = 1\,048\,575$) utilizando Python. La implementación usó la fórmula directa $\neg_L(x) = (2^L-1) - x$ para evitar truncamiento.

\begin{itemize}
    \item \textbf{Identidad verificada:} En todos los casos sin truncamiento, $\max(N_k) + \min(M_k) = 2^L - 1$.
    \item \textbf{Sincronía de extremos:} En el 100\% de los casos, el paso donde $N_k$ alcanza su máximo global coincide con el paso donde $M_k$ alcanza su mínimo global.
    \item \textbf{Teorema 1 confirmado:} Para toda trayectoria que converge a 1, $M_K = 2^L - 2$ en el mismo paso $K$.
\end{itemize}

% ==================== APÉNDICE B ====================
\section{Apéndice B: Análisis Computacional para $L=32$}
\label{app:L32}

Se extendió el marco a $L=32$ ($2^{32}-1 = 4\,294\,967\,295$) analizando $n = 1 \dots 10^5$. Ninguna trayectoria excedió $2^{32}$, por lo que no hubo truncamiento.

\subsection{Verificación Estadística de la Identidad Dual}

La correlación de Pearson entre $\max(N_k)$ y $\min(M_k)$ resultó ser exactamente $-1.0$ ($p \approx 0$, $R^2 = 1.0$). El ajuste lineal $\min(M) = a \cdot \max(N) + b$ arrojó $a = -1.000000$ y $b = 4\,294\,967\,295.00$, en concordancia exacta con la identidad teórica $M_k + N_k = 2^L - 1$.

\textbf{Nota:} Esta correlación perfecta es una identidad algebraica, no un hallazgo empírico. Su verificación confirma la correctitud de la implementación computacional.

\subsection{Sincronía de Extremos (Resultado Empírico No Trivial)}

En el \textbf{100\%} de las trayectorias analizadas ($10^5 / 10^5$), el índice $k$ donde $N_k$ alcanza su máximo global coincide exactamente con el índice donde $M_k$ alcanza su mínimo global:
\[
\arg\max_k N_k = \arg\min_k M_k \quad \text{(100\% de los casos verificados)}
\]

\subsection{Predictibilidad de la Longitud de Trayectoria (ML Exploratorio)}

Un modelo de Random Forest entrenado con características $\{n_0, \max(N_k), \min(M_k)\}$ logró un $R^2 \approx 0.48$ para predecir el número de pasos hasta 1. La variable $\log(\max(N_k))$ explicó el 46\% de la varianza, confirmando que la altura alcanzada es el predictor dominante de la duración de la trayectoria.

% ==================== APÉNDICE C ====================
\section{Apéndice C: Análisis Modular Empírico}
\label{app:modular}

Se analizaron $10^5$ trayectorias convergentes para investigar si la congruencia del número inicial $n_0$ predice propiedades estadísticas de la trayectoria.

\subsection{Sesgo en la Proporción de Pasos Pares}

La heurística estándar de Collatz predice una proporción de pasos pares tendiendo a $\frac{\log_2 3}{1+\log_2 3} \approx 0.6131$. En nuestra muestra, la proporción observada fue $\langle \text{even\_ratio} \rangle = 0.6845$, notablemente superior al valor teórico. Este sesgo positivo es consistente con correlaciones de corto alcance bien documentadas en la dinámica de Collatz \cite{lagarias}.

\subsection{Estructura Modular de la Longitud de Trayectoria}

Se aplicó el test no paramétrico de Kruskal-Wallis para evaluar si la distribución del número de pasos hasta 1 difiere según la clase de congruencia de $n_0$ (Tabla \ref{tab:kruskal}).

\begin{table}[h!]
\centering
\caption{Resultados del test de Kruskal-Wallis sobre la longitud de trayectoria por clase modular.}
\label{tab:kruskal}
\begin{tabular}{lcccl}
\toprule
Clasificación & $H$-estadístico & $p$-valor & Significativo & Interpretación \\
\midrule
$n_0 \bmod 2$  & 1337.72  & $7.17 \times 10^{-293}$ & Sí & La paridad inicial condiciona la estructura (esperado por definición de $C$). \\
$n_0 \bmod 3$  & 1.19     & $5.52 \times 10^{-1}$   & No & La congruencia módulo 3 no predice la longitud. El factor 3 se mezcla rápidamente. \\
$n_0 \bmod 4$  & 2579.59  & $\approx 0$             & Sí & Determina los dos primeros pasos de la iteración. \\
$n_0 \bmod 6$  & 1340.61  & $\approx 0$             & Sí & Hereda la estructura de la paridad. \\
$n_0 \bmod 8$  & 3844.40  & $\approx 0$             & Sí & Mayor resolución de la condición inicial binaria. \\
$n_0 \bmod 12$ & 2585.97  & $\approx 0$             & Sí & Combina efectos de módulo 3 y módulo 4. \\
\bottomrule
\end{tabular}
\end{table}

\subsection{Implicaciones para la Dualidad Binaria}

Bajo la transformación espejo $M_k = (2^L-1) - N_k$, la paridad se invierte: si $N_k$ es par, $M_k$ es impar, y viceversa (asumiendo $L \geq 1$). Por tanto, la alta significancia de $n_0 \bmod 2$ en la secuencia normal implica una estructura dual opuesta pero determinista en la secuencia espejo.

El hallazgo de que $n_0 \bmod 3$ \textbf{no} es significativo es particularmente relevante: sugiere que la complejidad de la trayectoria no reside en la congruencia respecto al coeficiente multiplicativo 3, sino en la interacción iterada entre la paridad y la división binaria. La dualidad por inversión de bits, al ser una operación intrínsecamente ligada a la base 2, está naturalmente alienada con la fuente real de la variabilidad dinámica.

\end{document}